\documentclass[11pt]{article}
\usepackage[margin=1in]{geometry}
\usepackage{amsmath,amssymb}
\usepackage{booktabs}
\usepackage[T1]{fontenc}
\usepackage{lmodern}
\usepackage[hidelinks]{hyperref}
\setlength{\parindent}{0pt}
\setlength{\parskip}{6pt}
\newcommand{\prew}{\mathrel{\preceq_w}}
\begin{document}
\begin{center}
{\Large\bfseries Response to the Second-Round Review}\\[6pt]
JBCB-1505, second major revision\\
Carlos Ramirez Ovalle
\end{center}

\textbf{Revised title:} Locating resolution-dependent behavioral correspondence
in qualitative DNA-damage response models.

We thank the reviewer for the close mathematical examination and for identifying
how the paper can better serve computational biologists. We audited definitions
and directional conclusions independently and substantially restructured the
presentation around DNA-damage response. The comment headings below summarize
the issues addressed, rather than claiming to quote the report verbatim.
Locations refer to the revised manuscript and its separate Supplementary
Validation document. The reasoning, principal numerical results, and limitations
needed to assess each response are included below; no local files or repository
documents need to be consulted to follow the answers.

\section*{R2.1. Example 1: whether both simulation directions hold}
\textbf{Response.} We respectfully retain the original direction after a
quantifier-level and independent computational audit. Definition 4 requires that,
for every related pair $(x,y)$ and transition $x\xrightarrow{a}x'$, some matching
successor $y'$ is chosen with $(x',y')$ itself in the simulation relation.
The matching condition must continue recursively from that particular pair.

To make the example explicit here, the late-choice system has initial state
$\ell_0$ and exactly the transitions
\[
\ell_0\xrightarrow{a}\ell_1,\qquad
\ell_1\xrightarrow{b}\ell_b,\qquad
\ell_1\xrightarrow{c}\ell_c.
\]
The early-choice system has initial state $e_0$ and exactly the transitions
\[
e_0\xrightarrow{a}e_b,\quad e_0\xrightarrow{a}e_c,\quad
e_b\xrightarrow{b}e'_b,\quad e_c\xrightarrow{c}e'_c.
\]
All remaining states are terminal, and neither system has silent transitions.
Here $X\prew Y$ means that $Y$ simulates $X$: every step of $X$ must have
a reply from $Y$. An explicit simulation from early to late is
\[
R_{E\to L}=\{(e_0,\ell_0),(e_b,\ell_1),(e_c,\ell_1),
(e'_b,\ell_b),(e'_c,\ell_c)\}.
\]
Both early $a$-moves match the sole late $a$-move; the selected $b$-only or
$c$-only continuation is matched by the late state offering both. Terminal
pairs impose no further obligations.

For the reverse direction, matching $\ell_0\xrightarrow{a}\ell_1$ requires
either $(\ell_1,e_b)$ or $(\ell_1,e_c)$ to be related. The former cannot match
$c$; the latter cannot match $b$. These exhaust the possible $a$-successors.
Different transitions from a related pair may choose different replies, but
the single initial late $a$-move cannot be matched by combining the two early
successor states. Thus early$\prew$late holds, late$\not\prew$early holds,
and weak bisimilarity fails. Both complete prefix-closed languages are
$\{\epsilon,a,ab,ac\}$: trace matching may independently choose the two early
branches for the two longer words.

\textbf{Changes made.} Example 1 in Appendix A now lists all transitions, the
witness relation, the two-case contradiction, and the distinction between
trace choices and recursive pair obligations. Section 3 refers readers to this
example. Definition 4 and the original direction are unchanged.

\textbf{Verification.} Production predicates, the existing game oracle, a new
game with independently constructed weak targets, and a direct supplied-relation
checker agree. The deletion certificate removes the two failed successor pairs
before the initial pair. Exact languages were checked by complete acyclic
enumeration and finite subset-product traversal without a depth cutoff.
mCRL2 independently confirms bisimulation and exact traces, and both simulation
directions using its strong preorder on these tau-free inputs. This does not
attribute arbitrary weak-preorder checking to mCRL2.
Every check returns early$\prew$late true, late$\prew$early false,
weak bisimilarity false, and exact trace equality true. The witness and
failed-successor argument above establish these conclusions independently
of access to any software output.

\textbf{Status.} Respectfully clarified; original mathematical direction
retained after audit. We agree that the previous exposition did not make the
successor-pair quantifiers sufficiently explicit.

\section*{R2.2. Proposition 2 and the roles of the strictness examples}
\textbf{Response.} We agree that each failed converse must have a clearly
identified witness. The hierarchy remains strong bisimilarity implies weak
bisimilarity implies mutual weak simulation implies exact weak-trace equality.
The forward implications follow from the transfer definitions and induction
along concrete finite paths.

\textbf{Changes made.} Proposition 2 and Table 5 in Appendix A distinguish:
\begin{enumerate}
\item Weak does not imply strong bisimilarity: $a.0$ versus $a.\tau.0$.
The silent move can be ignored weakly but has no direct strong match.
\item Mutual simulation does not imply weak bisimilarity:
$a.(b+c)+a.b$ versus $a.(b+c)$. This is the example highlighted by the reviewer,
already intended for the middle converse; its role is now explicit. Both
directions have witnesses, but a bisimulation must also match the extra
$b$-only successor against a state offering $c$.
\item Exact trace equality does not imply mutual weak simulation:
late choice versus early choice, with only early$\prew$late, as in Example 1.
\end{enumerate}
\textbf{Verification.} Explicit constructions, production code, independent
games, and exact languages give the following results. In process notation,
$a.P$ means an $a$-step followed by $P$, $+$ denotes a choice, $0$ is terminal,
and $\tau$ is silent.
\begin{center}
\small
\begin{tabular}{@{}lccccc@{}}
\toprule
Ordered pair $(X,Y)$ & Strong & Weak & $X\prew Y$ & $Y\prew X$ & Traces equal \\
\midrule
$(a.0,\,a.\tau.0)$ & No & Yes & Yes & Yes & Yes \\
$(a.(b+c)+a.b,\,a.(b+c))$ & No & No & Yes & Yes & Yes \\
$(\mathrm{late},\,\mathrm{early})$ & No & No & No & Yes & Yes \\
\bottomrule
\end{tabular}
\end{center}
Strong and Weak denote the two bisimilarities; the two simulation columns
are evaluated separately. mCRL2 confirms all three equivalence columns
and, for the two tau-free pairs, both simulation directions. No weak-preorder
claim is made for its check of the pair containing $\tau$.
\textbf{Status.} Addressed by clarification and independent audit.

\section*{R2.3. Systematic audit of directions and reported classes}
\textbf{Response.} We agree and extended the audit beyond Example 1.
In $L_1\prew L_2$, the right-hand system always simulates the left-hand system.

\textbf{Changes made.} A machine-readable audit records both directions for six
synthetic pairs, three GIM interfaces, five curated pairs, and 18
asynchronous/synchronous HPN pairs. The early/late benchmark retains the verified
class \emph{candidate simulated by reference}. The GIM table names both organisms;
RCD states that the animal model simulates the plant model, not conversely.
HPN retains its two Boolean direction fields without numerical ranking.
The notebook begins with GIM and includes the executable quantifier audit.

\textbf{Verification.} All 32 ordered model-pair records and three hierarchy
witnesses agree between production and independent calculations. All six
construction-level classes remain correct; 6/6 is retained only as software
verification. The expanded complete audit of 67,600 ordered pairs over 260
named one-/two-state LTSs found no simulation or weak-bisimulation disagreement
and no hierarchy or trace-inclusion violation. Reflexivity and transitivity
also hold on that universe. Specifically, the enumeration uses initial state
zero, alphabet $\{a,\tau\}$, and every edge subset on one or two named states:
$2^2+2^8=260$ systems and $260^2=67{,}600$ ordered pairs. It includes
41,080 weakly bisimilar pairs and 4,848 mutually simulated but non-bisimilar
pairs. There is no one-way simulation with exact trace equality in this small
universe; the larger Example 1 above supplies that witness. These finite
checks do not replace general proofs. Supplementary Validation S4 reports
the same scope and outcomes.

\textbf{Status.} Fully addressed within the declared models and audit universe.
No mathematical or directional correction was required.

\section*{R2.4. Fixed-point argument and implementation semantics}
\textbf{Response.} We agree that the fixed-point argument is sound and
rechecked it rather than changing its statement because the examples were
questioned. Enlarging a relation cannot remove matching witnesses, so each
operator is monotone. The greatest fixed point stays contained in the current
relation through every in-place deletion. At termination the retained pairs
satisfy the operator, and deleted pairs cannot become valid in a smaller
relation. The terminal relation is therefore the greatest fixed point.

\textbf{Changes made.} Theorem 1, the controlled-refinement and label-blindness
propositions, and their proofs remain in Appendix A. Definitions and the
executable algorithm remain in Section 3; the article is self-contained.
Silent moves continue to admit zero or more matching silent steps.

\textbf{Verification.} The in-place production algorithm was contrasted with
simultaneous least-losing-set rounds using independently built weak targets.
Both retain the same relations in the audits summarized in R2.3.
For completeness, let $\Phi$ be the matching operator and $R$ the current
relation, initially the full state-pair set. The invariant $\nu\Phi\subseteq R$
holds initially. If a pair in $\nu\Phi$ were deleted for lacking a witness
in $R$, its fixed-point witness in $\nu\Phi\subseteq R$ would contradict that
deletion. At termination, $R\subseteq\Phi(R)$, so $R\subseteq\nu\Phi$;
together the inclusions establish equality. Each successful deletion removes
one pair permanently, giving at most $n_1n_2$ deletions for systems with
$n_1$ and $n_2$ states. Trace inclusion follows by matching successive
observable steps and concatenating their weak replies. The independent audit
therefore supports, rather than substitutes for, the fixed-point and
path-induction arguments.
\textbf{Status.} Addressed by clarification and independent audit.

\section*{R2.5. A title better suited to computational biologists}
\textbf{Response.} We agree. The previous title foregrounded a formal framework.
The new title, \emph{Locating resolution-dependent behavioral correspondence in
qualitative DNA-damage response models}, names the biological modeling problem
without implying experimental validation or a universal law.

\textbf{Changes made.} The title, running head, abstract, and opening sections
are aligned with the GIM question.
\textbf{Verification.} Title, abstract, and first pages were checked together
for the models, readout question, principal result, and scope limitation.
\textbf{Status.} Addressed by substantial restructuring.

\section*{R2.6. Abstract emphasis on software metrics and inconclusive HPN}
\textbf{Response.} We agree and rewrote the abstract from scratch around
DNA-damage response, declared interfaces, constant structural similarity, and
the change in executable correspondence.

\textbf{Changes made.} The abstract omits 6/6, 42 decisions, 67,600 pairs, and
HPN $p=0.25$. It reports PN-GDDA 0.9976, weak bisimilarity at A/B, and failure
of both simulation directions at C, with an explicit model/interface limit.
Independent verification occupies one short supporting sentence.

\textbf{Verification.} The regenerated values are reproduced in the table in
R2.8 below. The revised abstract contains 170 words and foregrounds the
biological question, the A/B-to-C result, and its scope limitation.
\textbf{Status.} Addressed by substantial restructuring.

\section*{R2.7. Introduction and contribution statement}
\textbf{Response.} We agree that the biological comparison problem should
precede fixed-point operators.

\textbf{Changes made.} Section 1 is completely rewritten: executable model
comparison, concrete animal/plant DDR motivation, limits of structure and
traces, intuitive recursive matching, GIM results, and a bounded contribution.
Section 2 introduces models and literature-supported interfaces before
Section 3 defines the method. The contribution is locating correspondence
within declared biological readouts using established formal relations,
not proposing a new equivalence relation.

\textbf{Verification.} Software-validation counts are absent from the abstract
and early introduction; biological motivation precedes formal definitions.
\textbf{Status.} Addressed by substantial restructuring.

\section*{R2.8. Results organization and GIM as the principal result}
\textbf{Response.} We agree that GIM should lead the results rather than follow
an extended software-verification narrative.

\textbf{Changes made.} Section 4 now leads with fixed structure, coarse function,
pathway distinctions, terminal mechanisms, and the interpretation of changed
observability. Figure 2 places the biological abstractions and A/B/C outcomes
together. Supporting curated cases follow in Section 5; exploratory HPN follows
in Section 6.

\textbf{Verification.} Regeneration retains 13/14 states and PN-GDDA
0.9975845917160132. A/B retain weak bisimilarity and both simulations; C
retains neither simulation. The depth-six diagnostic changes from zero to
0.3529411764705882. The nets and initial markings are unchanged.
Here the two genome-integrity maintenance (GIM) models describe animal/human
and Arabidopsis DNA-damage responses. Interface A groups shared functions,
including a coarse terminal loss of proliferative capacity. B additionally
distinguishes damage channels, ATM/ATR pathways, and repair while retaining
a shared terminal label. C retains these pathway distinctions and separates
animal apoptosis from the encoded plant terminal mechanism. The latter
contains a combined differentiation/endoreduplication transition, not two
independent alternatives.
\begin{center}
\small
\begin{tabular}{@{}lccccc@{}}
\toprule
Interface & PN-GDDA & Weak & Animal$\prew$plant & Plant$\prew$animal & $d_6$ \\
\midrule
A: coarse & 0.9976 & Yes & Yes & Yes & 0 \\
B: pathway & 0.9976 & Yes & Yes & Yes & 0 \\
C: terminal & 0.9976 & No & No & No & 0.3529 \\
\bottomrule
\end{tabular}
\end{center}
Weak denotes weak bisimilarity. Strong bisimilarity fails at all three
interfaces. The diagnostic $d_6$ is the Jaccard distance between observable
trace sets truncated at six events; it is not an exact-trace decision.
The interpretation is a correspondence boundary between the declared B and C
readouts, conditional on these models. This is neither a new discovery of the
terminal mechanisms nor evidence of organism-level equivalence.
\textbf{Status.} Addressed by substantial restructuring.

\section*{R2.9. Verification should support, not dominate, the story}
\textbf{Response.} We agree and retained rigor while changing its location.

\textbf{Changes made.} Section 7 now provides a compact evidence summary in
Table 4, identifying each check, its independent basis, and the observed result:
6/6 constructed classes recovered; 42/42 external strong/weak decisions in
agreement; no disagreement on 32 model pairs and three hierarchy witnesses;
no discrepancy across 67,600 ordered small-system pairs; and PN-GDDA agreement
with Holmes to 12 decimal places on the stated control. The text defines the
scope of each check, distinguishes mathematical proof from finite testing,
and does not attribute general weak-preorder checking to mCRL2.
The full benchmark table/figure, comparator details, public-model controls, mCRL2
agreement, exhaustive finite checks, and scaling details are in Supplementary
Validation. Long proofs and quantifier examples are in Appendix A; full GIM
Petri nets remain in Appendix B. No underlying analysis or net diagram was
deleted.

\textbf{Verification.} Main and supplement compile separately with inline
bibliographies and resolved cross-references. Controls were regenerated, not
manually relabeled.
\textbf{Status.} Addressed by substantial restructuring.

\section*{R2.10. HPN-DREAM/CASPOTS remains exploratory}
\textbf{Response.} We agree. HPN does not validate predictive utility or the
GIM observational boundary.

\textbf{Changes made.} HPN is absent from the abstract and explicitly
exploratory in Section 6. Full methods, the direction-preserving plot, and
negative results remain in Supplementary Validation S3. The nine observations,
nominal direction categories, inconclusive diagnostic associations, and native
$p=0.25$ are retained. No more favorable statistical test was substituted.

\textbf{Verification.} Selection remains structure-only; hashes and directions
were checked. Asynchronous and synchronous analyses and exact permutations were
regenerated. Existing CASPOTS compatibility scores were retained, not
reinterpreted as unique time-series predictions; the CASPOTS optimizer was
not rerun for this revision. For transparency, the retained numerical evidence
is as follows. Of nine asynchronous pair--condition comparisons, three have
left-model-in-right-model simulation, three the reverse, and three neither;
none is weakly bisimilar. Median experimental profile RMSE in those nominal
categories is 0.2757, 0.1919, and 0.2071, respectively, with only three
observations per category. Synchronous updating preserves seven of nine classes.
Truncated-trace distance has Spearman $\rho=-0.596$ and exact one-sided
$p=0.556$; LTS-GDA distance has $\rho=0.669$ and $p=0.097$, using all
216 condition-stratified permutations. The CASPOTS native-context advantage
over mean cross-cell medoids is 0.00212 excess-RMSE units, not established
by the exact three-cell sign test ($p=0.25$). No ordinal ranking of formal
classes or predictive conclusion is supported by these results.
\textbf{Status.} Fully addressed in presentation and interpretation; small
sample size and inconclusive biological evidence remain substantive limits.

\section*{R2.11. Encoded mechanistic distinctions are not a discovery}
\textbf{Response.} We agree. Interfaces informed by known biology do not provide
a blinded discovery of the same mechanisms.

\textbf{Changes made.} Sections 2, 4.4--4.5, and 8.3 distinguish prior knowledge
of terminal mechanisms, the PN-GDDA structural description, the computed
survival/failure of recursive correspondence, and the unvalidated experimental
implication. The boundary is between B and C within three declared interfaces,
not a unique or globally minimal biological resolution. The combined plant
terminal transition is not described as two independent alternatives.

\textbf{Verification.} Only labels change between interfaces. Claim auditing
does not permit inference of organism-level equivalence, mechanistic identity,
prognosis, or experimental confirmation.
\textbf{Status.} Fully addressed by bounded interpretation. Experimental
validation remains future work, not an accomplished result.

\section*{R2.12. Figures should foreground the biological question}
\textbf{Response.} We agree.
\textbf{Changes made.} Figure 1 starts from the biological question, moves
through readout selection and executable comparison, and ends at a conditional
compatibility interpretation. Verification is a supporting layer. Figure 2
places animal/human and Arabidopsis functions above the A/B/C outcomes.
The benchmark and HPN plot become Supplementary Figures S1 and S2.
Detailed Petri nets remain main Figures 3 and 4 in Appendix B.
\textbf{Verification.} Figure 2 uses the regenerated values reproduced in R2.8.
Its functional categories are not presented as a required causal sequence.
The two complete Petri-net diagrams are preserved and remain legible in
the revised manuscript.
\textbf{Status.} Addressed by substantial restructuring.

\section*{R2.13. Discussion, conclusion, and accessible biological language}
\textbf{Response.} We agree that biological/computational meaning should lead
the interpretation.

\textbf{Changes made.} Section 8 asks what structural similarity tells us, what
behavioral comparison adds, and what interface refinement means biologically,
then treats execution assumptions, evidence limits, and computational scope.
The conclusion starts with the GIM A/B-to-C result. Biological sections explain
step/choice matching before formal names. Fixed-point correctness is no longer
the conclusion's opening claim.

\textbf{Verification.} The final narrative separates formal correctness,
implementation verification, support for executable models, and independent
biological evidence. Negative and inconclusive outcomes remain visible.
The manuscript and supplement follow the same interpretation summarized in
this response.
\textbf{Status.} Addressed by substantial restructuring.

\end{document}
